\chapter{Results and Analysis}
\label{chap:results}

This chapter presents the quantitative and qualitative results of the experiments outlined in the methodology. We evaluate the performance of both the \textbf{Fixed-Depth} and \textbf{Adaptive-Depth U-Net} models against a standard Bicubic interpolation baseline. All metrics are reported on the $256 \times 256$ cropped DIV2K validation set, using a 4-pixel border shave on the luminance (Y) channel.

%-----------------------------------------------------------
\section{Experiment 1: Fixed-Depth Model Performance}
\label{sec:results_fixed}

The first experiment evaluates the performance of a U-Net with a \textbf{fixed depth of 3 levels} across all scale factors. The objective is to establish a performance baseline for a non-adaptive architecture. The results are compared against the bicubic baseline in Table \ref{tab:results_fixed_depth}.

% Note: You must have \usepackage{booktabs} in your main .tex file
\begin{table}[h!]
    \centering
    \caption{Quantitative Results for Fixed-Depth (3 Levels) U-Net vs. Bicubic Baseline. The $\Delta$ column shows the improvement of our model over bicubic interpolation.}
    \label{tab:results_fixed_depth}
    \begin{tabular}{c c | c c c | c c c | c}
        \toprule
        \textbf{Scale} & \textbf{Fixed} & \multicolumn{3}{c}{\textbf{PSNR (dB)}} & \multicolumn{3}{c}{\textbf{SSIM}} & \textbf{MS-SSIM} \\
        \cmidrule(lr){3-5} \cmidrule(lr){6-8}
        \textbf{($s$)} & \textbf{Depth} & \textbf{Bicubic} & \textbf{Model} & \textbf{$\Delta$} & \textbf{Bicubic} & \textbf{Model} & \textbf{$\Delta$} & \textbf{(Model)} \\
        \midrule
        0.20 & 3 & & & & & & & \\
        0.30 & 3 & & & & & & & \\
        0.40 & 3 & & & & & & & \\
        0.50 & 3 & & & & & & & \\
        0.60 & 3 & & & & & & & \\
        0.70 & 3 & & & & & & & \\
        0.90 & 3 & & & & & & & \\
        \bottomrule
    \end{tabular}
\end{table}

\paragraph{Analysis of Fixed-Depth Results}
(TODO: Your analysis here. Discuss the trends from Table \ref{tab:results_fixed_depth}. Note the general improvement over bicubic and any potential performance bottlenecks at extreme scales.)

%-----------------------------------------------------------
\section{Experiment 2: Adaptive-Depth Model Performance}
\label{sec:results_adaptive}

The second experiment tests the core hypothesis: that adapting the network depth based on the scale factor improves performance. Deeper networks are used for higher-detail scales, while shallower networks are used for low-detail scales. The results are shown in Table \ref{tab:results_adaptive_depth}.

\begin{table}[h!]
    \centering
    \caption{Quantitative Results for \textbf{Adaptive-Depth U-Net} vs. Bicubic Baseline. The depth is now a function of the scale $s$.}
    \label{tab:results_adaptive_depth}
    \begin{tabular}{c c | c c c | c c c | c}
        \toprule
        \textbf{Scale} & \textbf{Adaptive} & \multicolumn{3}{c}{\textbf{PSNR (dB)}} & \multicolumn{3}{c}{\textbf{SSIM}} & \textbf{MS-SSIM} \\
        \cmidrule(lr){3-5} \cmidrule(lr){6-8}
        \textbf{($s$)} & \textbf{Depth} & \textbf{Bicubic} & \textbf{Model} & \textbf{$\Delta$} & \textbf{Bicubic} & \textbf{Model} & \textbf{$\Delta$} & \textbf{(Model)} \\
        \midrule
        0.20 & 1 & & & & & & & \\
        0.30 & 2 & & & & & & & \\
        0.40 & 3 & & & & & & & \\
        0.50 & 3 & & & & & & & \\
        0.60 & 4 & & & & & & & \\
        0.70 & 5 & & & & & & & \\
        0.80 & 6 & & & & & & & \\
        0.90 & 7 & & & & & & & \\
        \bottomrule
    \end{tabular}
\end{table}

\paragraph{Analysis of Adaptive-Depth Results}
(TODO: Your analysis here. Discuss the results from Table \ref{tab:results_adaptive_depth}, focusing on the $\Delta$ vs. Bicubic improvements. Set the stage for the direct comparison in the next section.)

%-----------------------------------------------------------
\section{Comparative Analysis: Fixed vs. Adaptive Depth}
\label{sec:results_comparison}

This section provides the direct comparison to answer the primary research question. Table \ref{tab:results_comparison} consolidates the model performance from both experiments to highlight the net benefit (or cost) of the adaptive-depth strategy.

\begin{table}[h!]
    \centering
    \caption{Direct Performance Comparison: Fixed-Depth (D=3) vs. Adaptive-Depth Model. The $\Delta$ columns represent the performance gain of the adaptive model over the fixed model.}
    \label{tab:results_comparison}
    \begin{tabular}{c | c c c | c c c}
        \toprule
        \textbf{Scale} & \multicolumn{3}{c}{\textbf{PSNR (dB)}} & \multicolumn{3}{c}{\textbf{SSIM}} \\
        \cmidrule(lr){2-4} \cmidrule(lr){5-7}
        \textbf{($s$)} & \textbf{Fixed (D=3)} & \textbf{Adaptive} & \textbf{$\Delta$ (Adaptive - Fixed)} & \textbf{Fixed (D=3)} & \textbf{Adaptive} & \textbf{$\Delta$ (Adaptive - Fixed)} \\
        \midrule
        0.20 & & & & & & \\
        0.30 & & & & & & \\
        0.40 & & & & & & \\
        0.50 & & & & & & \\
        0.60 & & & & & & \\
        0.70 & & & & & & \\
        0.80 & & & & & & \\
        0.90 & & & & & & \\
        \bottomrule
    \end{tabular}
\end{table}

\paragraph{Discussion of Comparative Results}
(TODO: This is your most important analysis. Use the $\Delta$ columns in Table \ref{tab:results_comparison} to prove or disprove your hypothesis.
\begin{itemize}
    \item At low scales ($s \le 0.4$): Does the adaptive model win?
    \item At high scales ($s \ge 0.6$): Does the adaptive model win?
    \item At crossover scales ($s=0.4, 0.5$): Is the performance similar, as expected?
\end{itemize}
)
%-----------------------------------------------------------
\section{Analysis of Training Dynamics}
\label{sec:results_training_dynamics}

Beyond the final checkpoint metrics, analyzing the training and validation loss curves provides insight into model stability, convergence speed, and overfitting.

\subsection{Experiment 1: Fixed-Depth Training}
Figure \ref{fig:train_fixed} plots the validation PSNR (or loss) over training epochs for the fixed-depth (D=3) model at different scales.

\begin{figure}[h!]
    \centering
    % TODO: Plot your training/validation curves
    % You will need to export these from TensorBoard
    % \includegraphics[width=0.49\linewidth]{path/to/fixed_depth_s03_loss.png}
    % \includegraphics[width=0.49\linewidth]{path/to/fixed_depth_s07_loss.png}
    
    \framebox{\parbox{0.9\textwidth}{\centering
        \vspace{5cm}
        \textbf{Graph Placeholder: Fixed-Depth (D=3) Training Curves} \\
        \small Plot validation PSNR (Y-axis) vs. Epoch (X-axis).
        \newline
        Use different colors for each scale (e.g., $s=0.3$, $s=0.5$, $s=0.7$, $s=0.9$).
        \vspace{5cm}
    }}
    \caption{Validation PSNR curves for the \textbf{Fixed-Depth (D=3)} model across various scales. (You may also want a second plot showing the train/validation loss gap for a specific scale).}
    \label{fig:train_fixed}
\end{figure}

\paragraph{Analysis of Fixed-Depth Training}
(TODO: Your analysis here.
\begin{itemize}
    \item Do higher scales (e.g., 0.9) converge faster or slower?
    \item Do they plateau at a higher or lower PSNR?
    \item Do the low scales (e.g., 0.2, 0.3) show signs of overfitting (a large gap between training and validation loss) because the model is too complex for the simple task?)
\end{itemize}
)

\subsection{Experiment 2: Adaptive-Depth Training}
Figure \ref{fig:train_adaptive} shows the same plot, but for the adaptive-depth models. This compares, for example, the convergence of the $s=0.3$ (D=2) model against the $s=0.7$ (D=5) model.

\begin{figure}[h!]
    \centering
    % TODO: Plot your training/validation curves
    % \includegraphics[width=0.49\linewidth]{path/to/adaptive_depth_s03_loss.png}
    % \includegraphics[width=0.49\linewidth]{path/to/adaptive_depth_s07_loss.png}
    
    \framebox{\parbox{0.9\textwidth}{\centering
        \vspace{5cm}
        \textbf{Graph Placeholder: Adaptive-Depth Training Curves} \\
        \small Plot validation PSNR (Y-axis) vs. Epoch (X-axis).
        \newline
        Use different colors for each scale/depth pair (e.g., $s=0.3/D=2$, $s=0.5/D=3$, $s=0.7/D=5$).
        \vspace{5cm}
    }}
    \caption{Validation PSNR curves for the \textbf{Adaptive-Depth} models across various scales.}
    \label{fig:train_adaptive}
\end{figure}

\paragraph{Analysis of Adaptive-Depth Training}
(TODO: Your analysis here.
\begin{itemize}
    \item How do these curves compare to the fixed-depth ones in Figure \ref{fig:train_fixed}?
    \item Does the $s=0.3$ (D=2) model show a smaller validation gap than the $s=0.3$ (D=3) model from the first experiment (suggesting less overfitting)?
    \item Does the $s=0.7$ (D=5) model converge to a higher validation PSNR than the $s=0.7$ (D=3) model?
    \item Are the training curves generally more stable?)
\end{itemize}
)
%-----------------------------------------------------------
\section{Qualitative Results (Visual Comparison)}
\label{sec:results_qualitative}

Quantitative metrics provide an objective measure of error, but qualitative (visual) inspection is necessary to assess perceptual quality, such as texture reconstruction and artifact suppression. Figures \ref{fig:qual_low_scale} and \ref{fig:qual_high_scale} show visual comparisons on selected patches from the validation set.

% You will need to \usepackage{graphicx} in your main .tex file

\begin{figure}[h!]
    \centering
    % TODO: Create a 1x5 grid of your output images
    % \includegraphics[width=0.19\linewidth]{path/to/img_s03_lr_input.png}
    % \includegraphics[width=0.19\linewidth]{path/to/img_s03_bicubic.png}
    % \includegraphics[width=0.19\linewidth]{path/to/img_s03_fixed.png}
    % \includegraphics[width=0.19\linewidth]{path/to/img_s03_adaptive.png}
    % \includegraphics[width=0.19\linewidth]{path/to/img_s03_hr.png}
    
    % Placeholder for now
    \framebox{\parbox{0.9\textwidth}{\centering
        \vspace{5cm}
        \textbf{Image Placeholder: Low-Scale (e.g., $s=0.3$)} \\
        \small (a) LR Input \quad (b) Bicubic \quad (c) Fixed-Depth (D=3) \quad (d) Adaptive-Depth (D=2) \quad (e) Ground Truth
        \vspace{5cm}
    }}
    
    \caption{Visual comparison on a low-detail scale ($s=0.3$).}
    \label{fig:qual_low_scale}
\end{figure}

\begin{figure}[h!]
    \centering
    % TODO: Create a 1x5 grid of your output images
    % \includegraphics[width=0.19\linewidth]{path/to/img_s07_lr_input.png}
    % \includegraphics[width=0.19\linewidth]{path/to/img_s07_bicubic.png}
    % \includegraphics[width=0.19\linewidth]{path/to/img_s07_fixed.png}
    % \includegraphics[width=0.19\linewidth]{path/to/img_s07_adaptive.png}
    % \includegraphics[width=0.19\linewidth]{path/to/img_s07_hr.png}
    
    \framebox{\parbox{0.9\textwidth}{\centering
        \vspace{5cm}
        \textbf{Image Placeholder: High-Scale (e.g., $s=0.7$)} \\
        \small (a) LR Input \quad (b) Bicubic \quad (c) Fixed-Depth (D=3) \quad (d) Adaptive-Depth (D=5) \quad (e) Ground Truth
        \vspace{5cm}
    }}
    
    \caption{Visual comparison on a high-detail scale ($s=0.7$).}
    \label{fig:qual_high_scale}
\end{figure}

\paragraph{Analysis of Qualitative Results}
(TODO: Discuss the visual evidence from your figures. Point to specific areas (e.g., "in Figure \ref{fig:qual_high_scale} (d), the textures..."). Explain *how* the visual results support the quantitative findings in your tables.)